\section{Notation and Gaussian completion}\label{sec:setup}

Fix a finite $\beta>1$. Constants may depend on this fixed parameter;
when that dependence matters we write $C_\beta$ or $O_\beta(\cdot)$.
All logarithms are natural. For probability measures $\mu\ll\nu$, set
\[
 {\rm KL}(\mu\Vert\nu)=\int\log\frac{d\mu}{d\nu}\,d\mu.
\]
On the spin cube, total variation is one half of the $\ell^1$ distance.
For vectors in $\mathbb R^n$ write
$\langle v,w\rangle_n=v\cdot w/n$ and
$|v|_n^2=\langle v,v\rangle_n$. The notation $\|\cdot\|$ refers to
the ordinary Euclidean norm, or to the operator norm of a matrix.

The orthogonal queries will use a Gaussian orthogonal ensemble (GOE)
matrix. Complete the interaction matrix by independent diagonal
entries: $W=W(G)$ is symmetric, its upper off-diagonal entries are
$J_{ij}/\sqrt n$, and its diagonal entries have law $N(0,2/n)$.
The vector $G$ consists of the independent standard Gaussian drivers
of all these entries, and $\gamma$ denotes its law. Define
\[
 \widehat H_G(x)=\frac12x^{\mathsf T}W(G)x,
 \qquad F(G)=\log\sum_{x\in\{-1,1\}^n}e^{\beta\widehat H_G(x)},
 \qquad \pi_G(x)=e^{\beta\widehat H_G(x)-F(G)}.
\]
Since $x_i^2=1$, the added diagonal contributes the same centered
random constant to every spin energy. It changes neither $\pi_G$
nor the heat-bath transition rates, and its contribution to
$\mathbb E F(G)$ vanishes.

The joint law to which proposals will be compared is
\[
 P(dG,x)=\gamma(dG)\pi_G(x).
\]
Thus $P$ samples the ordinary disorder first and then a Gibbs spin.
A proposal law $Q_l$ will also have disorder marginal $\gamma$, but
its conditional spin law will be constructed by matrix queries.
Keeping these joint laws distinct is necessary: entropy on the joint
space first controls proposals, whereas the desired conclusion concerns
a transition law for one fixed disorder and one fixed initial spin.
